% ============================================================
% CAPÍTULO 3 — INTELIGÊNCIA ARTIFICIAL E OS DESAFIOS 
%              À PERSONALIDADE HUMANA
% ============================================================

\chapter{Inteligência Artificial e os Desafios à Personalidade Humana}
\label{cap:ia-desafios}

\section{Introdução ao Capítulo}
\label{sec:intro-cap3}

O avanço da inteligência artificial nas últimas décadas tem produzido transformações profundas em praticamente todas as esferas da vida social. Sistemas algorítmicos já não se limitam a tarefas repetitivas ou de processamento de dados: hoje participam de processos decisórios que afetam diretamente a liberdade, a privacidade, a igualdade e a identidade das pessoas. Essa penetração capilar da IA na vida cotidiana suscita questões jurídicas de primeira grandeza, que desafiam categorias tradicionais do direito civil, constitucional e dos direitos humanos.

A \MH{} dedica especial atenção a esses desafios, particularmente em seu Capítulo III, parágrafos 90 a 130, onde oferece uma análise dos impactos da IA sobre a pessoa humana à luz dos princípios antropológicos e éticos examinados no capítulo anterior. A Encíclica não se limita a diagnosticar riscos, mas propõe critérios para o discernimento ético e jurídico das tecnologias emergentes.

Este capítulo examina, em diálogo com a \MH{} e com a literatura jurídica especializada, os principais desafios que a IA impõe à personalidade humana. O percurso é o seguinte: parte-se de uma caracterização da IA e de seus fundamentos técnicos (\ref{sec:fundamentos-ia}), para então analisar os desafios específicos à autonomia (\ref{sec:automacao}), à igualdade (\ref{sec:discriminacao}), à privacidade (\ref{sec:vigilancia}), à identidade pessoal (\ref{sec:identidade}) e à agência moral (\ref{sec:agencia-moral}).

\section{Fundamentos Técnicos e Conceituais da IA}
\label{sec:fundamentos-ia}

\subsection{Definição e Tipologias}
\label{subsec:definicao-ia}

A definição de inteligência artificial é objeto de controvérsia na literatura técnica e jurídica. Para os fins deste trabalho, adota-se o conceito formulado pelo Grupo de Especialistas de Alto Nível em IA da União Europeia: sistemas de software (e eventualmente hardware) projetados por seres humanos que, dado um objetivo complexo, agem na dimensão física ou digital percebendo seu ambiente, interpretando os dados coletados, raciocinando sobre o conhecimento derivado desses dados e decidindo as melhores ações a serem tomadas para atingir o objetivo \cite{HighLevelExpertGroup2019EA}.

A Encíclica \MH{} oferece uma definição complementar, de caráter antropológico: a IA é apresentada como ``um conjunto de tecnologias que simulam capacidades humanas de aprendizado, raciocínio e decisão, mas que operam em um registro fundamentalmente diverso da inteligência humana, por carecerem de consciência, liberdade e intencionalidade moral'' \cite{LeaoXIV2026MH}.

Essa distinção é crucial. Embora os sistemas de IA possam executar tarefas cognitivas com eficiência superior à humana em domínios específicos, eles não possuem as qualidades que caracterizam a pessoa: subjetividade, autoconsciência, capacidade de atribuir sentido e valor às suas ações, e abertura ao transcendente. A \MH{} adverte contra o que chama de ``equívoco categorial'': tratar a inteligência artificial como se fosse um tipo de inteligência humana, quando na verdade se trata de fenômenos ontologicamente distintos.

\subsection{A Arquitetura dos Sistemas Algorítmicos}
\label{subsec:arquitetura-algoritmica}

A compreensão dos desafios jurídicos da IA requer um conhecimento mínimo de sua arquitetura técnica. Os sistemas contemporâneos de IA, especialmente aqueles baseados em aprendizado de máquina (\textit{machine learning}), operam em três etapas fundamentais: coleta e processamento de dados, treinamento de modelos e inferência preditiva.

Na \textbf{etapa de coleta}, grandes volumes de dados pessoais e impessoais são agregados a partir de múltiplas fontes — navegação na internet, sensores, transações, interações em redes sociais, registros públicos. Esses dados são então processados e estruturados para alimentar os algoritmos de treinamento.

Na \textbf{etapa de treinamento}, os modelos algorítmicos aprendem padrões estatísticos a partir dos dados fornecidos. O aprendizado pode ser supervisionado (com exemplos rotulados), não supervisionado (sem rótulos, identificando padrões latentes) ou por reforço (aprendendo por tentativa e erro com base em recompensas).

Na \textbf{etapa de inferência}, o modelo treinado é aplicado a novos dados para produzir predições, classificações ou recomendações. É nessa etapa que os sistemas de IA impactam diretamente a vida das pessoas: decidindo quem recebe um empréstimo, quem é chamado para uma entrevista de emprego, quem tem acesso a um seguro de saúde, quanto tempo alguém permanece detido sob suspeita criminal.

A \MH{} demonstra preocupação especial com a opacidade desse processo. A Encíclica observa que ``a complexidade dos algoritmos frequentemente supera a capacidade de compreensão de seus próprios criadores, criando uma assimetria de poder entre aqueles que controlam os sistemas e aqueles que são por eles afetados'' \cite{LeaoXIV2026MH}. Essa constatação tem implicações jurídicas diretas, pois a opacidade algorítmica compromete o exercício de direitos fundamentais como o contraditório e a ampla defesa.

\section{Automação e Autonomia Decisória}
\label{sec:automacao}

\subsection{A Delegação de Decisões a Sistemas Algorítmicos}
\label{subsec:delegacao-decisoria}

Um dos fenômenos mais significativos da era digital é a crescente delegação de decisões — públicas e privadas — a sistemas algorítmicos. No setor público, algoritmos são utilizados para a triagem de candidatos a políticas sociais, para a predição de reincidência criminal, para a alocação de recursos em saúde e educação, e até para a fiscalização tributária. No setor privado, a IA decide sobre concessão de crédito, precificação de seguros, recrutamento e seleção de pessoal, e recomendação de conteúdos.

Essa delegação levanta questões fundamentais sobre a autonomia humana. A \MH{} adverte que ``a tentação de confiar decisões importantes a máquinas, sob o pretexto de maior eficiência ou imparcialidade, pode levar a uma progressiva atrofia da capacidade humana de julgar e deliberar'' \cite{LeaoXIV2026MH}.

O direito brasileiro já reconhece a gravidade desse problema. O art. 20 da Lei Geral de Proteção de Dados Pessoais (Lei nº 13.709/2018) assegura ao titular o direito de solicitar revisão de decisões tomadas exclusivamente com base em tratamento automatizado de dados pessoais que afetem seus interesses. O §1º do mesmo artigo determina que a revisão deve ser realizada por pessoa natural, estabelecendo uma garantia mínima contra a desumanização dos processos decisórios \cite{Brasil2018LGPD}.

Contudo, a eficácia desse dispositivo tem sido questionada. Estudos empíricos mostram que, na prática, as solicitações de revisão raramente resultam em alteração da decisão automatizada, seja porque os mecanismos de recurso são obscuros ou inacessíveis, seja porque o revisor humano tende a confirmar a decisão algorítmica por efeito de ancoragem cognitiva \cite{Binns2018ITG, Edwards2017PC}.

\subsection{O Direito à Explicação}
\label{subsec:direito-explicacao}

O direito à explicação das decisões automatizadas emerge como uma das garantias fundamentais para a proteção da personalidade na era da IA. Esse direito encontra fundamento no art. 20, §2º, da LGPD, que obriga o controlador a fornecer informações claras e adequadas sobre os critérios e procedimentos utilizados na decisão automatizada \cite{Brasil2018LGPD}.

No direito europeu, o art. 22 do Regulamento Geral de Proteção de Dados (GDPR) estabelece regime semelhante, proibindo decisões baseadas exclusivamente em tratamento automatizado que produzam efeitos jurídicos ou afetem significativamente o titular dos dados, salvo nas exceções previstas \cite{UE2016GDPR}. O Grupo de Trabalho do Art. 29 (predecessor do Comitê Europeu de Proteção de Dados) já se manifestou no sentido de que o direito à explicação inclui informações sobre a lógica subjacente ao funcionamento do sistema, sua importância e as consequências previstas do tratamento.

A \MH{} endossa essa perspectiva ao afirmar que ``a transparência dos sistemas algorítmicos não é apenas uma exigência técnica, mas um imperativo ético que decorre do reconhecimento da pessoa como sujeito racional, que tem direito a compreender as decisões que afetam sua vida'' \cite{LeaoXIV2026MH}. A Encíclica vincula, assim, o direito à explicação à própria dignidade da pessoa humana: tratar alguém como sujeito de direitos implica reconhecer sua capacidade de compreender e questionar as decisões que lhe dizem respeito.

\subsection{O Desafio dos Modelos de Caixa-Preta}
\label{subsec:caixa-preta}

Um obstáculo técnico significativo ao direito à explicação é a natureza de ``caixa-preta'' de muitos sistemas de IA, especialmente os baseados em redes neurais profundas. Esses modelos atingem alto desempenho preditivo precisamente porque sua complexidade interna é tão grande que nem mesmo seus criadores conseguem explicar, em linguagem compreensível por seres humanos, como uma determinada decisão foi alcançada.

Esse fenômeno, conhecido como o trade-off entre acurácia e explicabilidade, coloca um dilema para o direito: se a proteção efetiva dos direitos fundamentais exige decisões explicáveis, e se os sistemas mais precisos são justamente os menos explicáveis, como conciliar esses dois imperativos?

A literatura jurídica tem proposto diferentes soluções para esse dilema. Uma primeira abordagem defende a proibição do uso de modelos de caixa-preta em decisões de alto risco, obrigando que sistemas utilizados em contextos penal, sanitário, creditício e laboratorial sejam necessariamente interpretáveis. Uma segunda abordagem aposta no desenvolvimento de técnicas de IA explicável (XAI — \textit{eXplainable Artificial Intelligence}) que permitam gerar explicações \textit{post hoc} mesmo para modelos complexos. Uma terceira abordagem, mais moderada, admite a utilização de modelos complexos desde que acompanhados de salvaguardas processuais robustas, como auditoria independente, testes de viés e supervisão humana efetiva \cite{Goodman2017EU, Selbst2018TIR, DoshiVelez2017MAR}.

A \MH{}, sem entrar no mérito técnico do debate, oferece critérios éticos para sua resolução: ``quando houver conflito entre a eficiência de um sistema e a proteção da dignidade humana, esta deve sempre prevalecer'' \cite{LeaoXIV2026MH}. Esse critério, se acolhido pelo ordenamento jurídico, implica que a explicabilidade deve ser considerada um requisito mínimo para a legitimidade do uso de IA em decisões que afetem direitos fundamentais.

\section{Discriminação Algorítmica}
\label{sec:discriminacao}

\subsection{Os Mecanismos da Discriminação Algorítmica}
\label{subsec:mecanismos-discriminacao}

A discriminação algorítmica ocorre quando sistemas de IA produzem resultados que desfavorecem sistematicamente determinados grupos ou indivíduos com base em características protegidas como raça, gênero, idade, orientação sexual, religião ou origem geográfica. Esse fenômeno pode decorrer de múltiplos fatores, que a literatura tem classificado em três categorias principais \cite{Noble2018AFS, Barocas2016BFF}.

A \textbf{discriminação por dados de treinamento} ocorre quando o conjunto de dados utilizado para treinar o modelo reflete preconceitos históricos existentes na sociedade. Se um algoritmo de recrutamento for treinado com dados históricos de contratação de uma empresa que discriminou mulheres, ele tenderá a reproduzir esse viés, excluindo candidatas do sexo feminino independentemente de sua qualificação.

A \textbf{discriminação por codificação} ocorre quando os próprios critérios e variáveis escolhidos para o modelo refletem escolhas valorativas que prejudicam determinados grupos. Por exemplo, um algoritmo de predição de reincidência criminal que utilize o histórico de prisões como variável preditiva tende a desfavorecer comunidades negras, que historicamente são alvo de policiamento mais intensivo e seletivo.

A \textbf{discriminação por \textit{proxy}} ocorre quando o algoritmo utiliza variáveis aparentemente neutras que funcionam como substitutas (\textit{proxies}) de características protegidas. O CEP de residência, por exemplo, pode funcionar como \textit{proxy} de raça ou classe social em muitas cidades brasileiras, produzindo discriminação indireta mesmo quando o modelo não utiliza explicitamente essas variáveis proibidas.

\subsection{A Discriminação Algorítmica na \MH{}}
\label{subsec:discriminacao-mh}

A \MH{} aborda a discriminação algorítmica como uma violação do princípio da igualdade e da justiça social. A Encíclica denuncia que a IA ``pode perpetuar e até ampliar as desigualdades existentes, criando um círculo vicioso no qual os já marginalizados são ainda mais prejudicados por sistemas que se apresentam como neutros e objetivos'' \cite{LeaoXIV2026MH}.

O documento estabelece uma conexão direta entre discriminação algorítmica e exclusão digital. Aqueles que não têm acesso à tecnologia, ou que não compreendem seu funcionamento, são duplamente vulneráveis: ficam excluídos dos benefícios da IA ao mesmo tempo em que podem ser prejudicados por suas decisões sem ter consciência disso ou meios de defesa. Essa constatação fundamenta a exigência de que a regulação da IA seja acompanhada de políticas de inclusão digital e de letramento algorítmico.

\subsection{Enfrentamento Jurídico da Discriminação Algorítmica no Brasil}
\label{subsec:enfrentamento-juridico-brasil}

O ordenamento jurídico brasileiro dispõe de instrumentos para o enfrentamento da discriminação algorítmica, embora ainda careça de um marco regulatório específico para a IA. O princípio da igualdade, inscrito no caput do art. 5º da Constituição Federal, combinado com a proibição de discriminação em razão de origem, raça, sexo, cor, idade e quaisquer outras formas de discriminação, constitui o fundamento primeiro para a tutela contra a discriminação algorítmica \cite{Brasil1988CF}.

A Lei de Igualdade Racial (Lei nº 12.288/2010), o Estatuto da Pessoa com Deficiência (Lei nº 13.146/2015) e a Lei de Combate à Discriminação (Lei nº 7.716/1989) oferecem instrumentos específicos de proteção. A LGPD, por sua vez, proíbe o tratamento de dados pessoais sensíveis para fins discriminatórios (art. 11) e determina que a revisão de decisões automatizadas deve levar em conta a não discriminação (art. 20, §2º) \cite{Brasil2018LGPD}.

No entanto, a efetividade desses instrumentos depende de adaptações interpretativas para o contexto algorítmico. A discriminação algorítmica frequentemente opera de forma indireta e difusa, tornando difícil a demonstração do nexo causal entre o funcionamento do sistema e o resultado discriminatório. A doutrina tem proposto, nesse sentido, a inversão do ônus da prova em casos de suspeita de discriminação algorítmica, bem como a instituição de auditorias obrigatórias de viés para sistemas de IA de alto risco \cite{Wachter2017TSP, Selbst2018TIR}.

\section{Vigilância, Controle e Privacidade}
\label{sec:vigilancia}

\subsection{O Capitalismo de Vigilância}
\label{subsec:capitalismo-vigilancia}

O conceito de capitalismo de vigilância, formulado por Shoshana Zuboff, descreve o modelo econômico emergente no qual a experiência humana é convertida em matéria-prima para a predição e o controle comportamental. Nesse modelo, os dados pessoais não são meramente um subproduto da atividade digital, mas o próprio objeto do processo produtivo: empresas de tecnologia extraem, analisam e monetizam dados comportamentais em escala massiva, utilizando algoritmos de IA para predizer e influenciar o comportamento futuro dos usuários \cite{Zuboff2019SC}.

A \MH{} dialoga com essa crítica ao denunciar que ``a coleta massiva de dados pessoais e sua utilização por sistemas algorítmicos cria uma assimetria de poder sem precedentes, na qual as pessoas são observadas, classificadas e manipuladas sem que tenham pleno conhecimento ou controle sobre esse processo'' \cite{LeaoXIV2026MH}.

O capitalismo de vigilância desafia a proteção jurídica da personalidade em pelo menos três dimensões. Em primeiro lugar, compromete a privacidade informacional, na medida em que dados pessoais são coletados e processados sem consentimento informado ou para finalidades diversas daquelas que motivaram a coleta. Em segundo lugar, afeta a autodeterminação informativa — o direito de controlar as próprias informações e decidir sobre sua circulação. Em terceiro lugar, ameaça a liberdade de escolha, pois a manipulação comportamental baseada em perfis algorítmicos pode induzir decisões que não refletem genuinamente a vontade do indivíduo.

\subsection{Vigilância Digital e Estado de Direito}
\label{subsec:vigilancia-estado-direito}

A utilização de sistemas de IA para fins de vigilância estatal suscita questões ainda mais graves do ponto de vista do Estado de Direito. Sistemas de reconhecimento facial em espaços públicos, análise preditiva de comportamento criminal, monitoramento massivo de comunicações e perfilamento ideológico são tecnologias que, se não submetidas a controles rigorosos, podem converter o Estado em Leviatã algorítmico.

O Supremo Tribunal Federal tem se manifestado crescentemente sobre essas questões. No julgamento da ADPF nº 695 (caso do reconhecimento facial no Metrô de São Paulo), o STF determinou que a utilização de tecnologias de reconhecimento facial pelo poder público exige lei específica, proporcionalidade, transparência e mecanismos efetivos de controle \cite{STF2022ADPF695}. Esse entendimento alinha-se à jurisprudência da Corte Europeia de Direitos Humanos, que no caso \textit{S. e Marper v. Reino Unido} (2008) já estabelecia limites rigorosos à retenção de dados biométricos pelo Estado.

A \MH{} contribui com essa discussão ao afirmar que ``a segurança pública não pode ser invocada como pretexto para a eliminação da privacidade e da liberdade'' \cite{LeaoXIV2026MH}. A Encíclica rejeita o falso dilema entre segurança e liberdade, sustentando que uma sociedade que sacrifica a privacidade em nome da segurança acaba por comprometer ambas.

\section{Identidade Pessoal e Subjetividade}
\label{sec:identidade}

\subsection{A Fragmentação da Identidade na Era Digital}
\label{subsec:fragmentacao-identidade}

A IA desafia também a integridade da identidade pessoal. Diferentemente das sociedades pré-digitais, onde a identidade de uma pessoa era construída principalmente por meio de narrativas biográficas e relações intersubjetivas, na era digital a identidade é cada vez mais definida por perfis algorítmicos construídos a partir de dados fragmentários — históricos de navegação, transações, localizações geográficas, interações em redes sociais.

Esses perfis não são neutros: eles refletem os interesses e as visões de mundo das empresas que os constroem. O perfil algorítmico de uma pessoa pode ser radicalmente diferente de sua autoimagem, e no entanto é esse perfil que determina seu acesso a crédito, emprego, seguros e serviços públicos. Cria-se, assim, uma tensão entre a identidade narrada (a história que a pessoa conta sobre si mesma) e a identidade algorítmica (o perfil que os sistemas atribuem à pessoa).

A \MH{} reconhece essa tensão ao afirmar que ``a pessoa humana não pode ser reduzida a um perfil ou a um conjunto de dados; ela é sempre mais do que aquilo que as máquinas podem capturar sobre ela'' \cite{LeaoXIV2026MH}. Esse reconhecimento fundamenta, no plano jurídico, o direito de não ser submetido a decisões baseadas exclusivamente em perfis algorítmicos, bem como o direito de contestar e corrigir as informações que compõem o perfil.

\subsection{A Voz Sintética e a Aparência Digital}
\label{subsec:voz-sintetica}

Tecnologias de IA generativa tornaram possível a criação de vozes sintéticas, imagens e vídeos hiper-realistas (deepfakes) que imitam pessoas reais com precisão perturbadora. Essas tecnologias abrem possibilidades criativas, mas também criam riscos graves para a identidade pessoal.

A clonagem não autorizada da voz ou da imagem de uma pessoa viola seus direitos de personalidade, especialmente o direito à própria imagem (art. 5º, V e X, da CF/88; arts. 12 a 20 do Código Civil). A \MH{} adverte que ``a criação de representações artificiais de pessoas reais, sem seu consentimento e especialmente com o intuito de enganar ou manipular, constitui uma forma de violência simbólica contra a pessoa, que perde o controle sobre sua própria representação'' \cite{LeaoXIV2026MH}.

Além disso, os deepfakes representam uma ameaça à confiança social necessária para o funcionamento das instituições democráticas. Vídeos fraudulentos de autoridades públicas podem ser utilizados para disseminar desinformação, influenciar eleições e desestabilizar governos. O direito à identidade pessoal, nesse contexto, adquire uma dimensão coletiva: proteger a integridade da representação de cada pessoa é também proteger a possibilidade de uma comunicação social confiável \cite{Chesney2019DF, Pariser2022THP}.

\section{Transumanismo, Pós-Humanismo e os Limites do Humano}
\label{sec:agencia-moral}

\subsection{O Projeto Transumanista}
\label{subsec:projeto-transumanista}

O transumanismo é um movimento intelectual e tecnológico que defende o uso da tecnologia para aprimorar as capacidades humanas para além de seus limites biológicos naturais, visando em última instância a superação da própria condição humana. Seus proponentes advogam a integração cada vez mais estreita entre seres humanos e máquinas, a extensão radical da longevidade, o aumento da inteligência por meios artificiais e, em suas versões mais radicais, o upload da consciência humana para suportes digitais \cite{Bostrom2005TH, Kurzweil2005THS}.

A \MH{} dedica atenção considerável ao transumanismo, identificando-o como uma das expressões mais extremas do paradigma tecnocrático. A Encíclica adverte que ``o desejo de superar os limites da condição humana, se não for acompanhado por uma reflexão profunda sobre o sentido da existência e sobre o bem integral da pessoa, pode levar a uma forma de desumanização mais sutil e perigosa do que as abertamente desumanizadoras'' \cite{LeaoXIV2026MH}.

A crítica da \MH{} ao transumanismo não se confunde com uma rejeição do progresso tecnológico. A Encíclica reconhece que intervenções tecnológicas no corpo e na mente humanos podem ser legítimas quando orientadas para a cura de patologias ou para a melhoria da qualidade de vida. A objeção recai sobre a pretensão de transcender a natureza humana como tal, substituindo a pessoa por uma entidade pós-humana que teria perdido precisamente aquilo que a torna humana: a corporeidade, a finitude, a historicidade e a capacidade de sofrer.

\subsection{Implicações Jurídicas do Transumanismo}
\label{subsec:implicacoes-juridicas}

O transumanismo levanta questões jurídicas profundas que desafiam categorias fundamentais do direito. Se a pessoa humana puder ser radicalmente modificada por intervenções tecnológicas, o que acontece com a noção de dignidade humana? Se a consciência puder ser transferida para suportes digitais, quem é o titular dos direitos de personalidade? Se a longevidade for estendida por séculos, como se reorganizam as relações contratuais, familiares e sucessórias?

A \MH{} não oferece respostas prontas a essas questões, mas estabelece um princípio orientador: ``a tecnologia deve servir à pessoa humana, e não o contrário; qualquer desenvolvimento tecnológico que reduza ou comprometa a humanidade do ser humano deve ser rejeitado'' \cite{LeaoXIV2026MH}. Esse princípio, se incorporado ao discurso jurídico, fundamenta a proibição de determinadas intervenções tecnológicas no corpo e na mente humanos, bem como a exigência de que o desenvolvimento tecnológico seja orientado por uma avaliação cuidadosa de seus impactos antropológicos.

\section{Síntese do Capítulo}
\label{sec:sumario-cap3}

Os desafios examinados neste capítulo demonstram que a IA não é neutra do ponto de vista antropológico e jurídico. Cada inovação tecnológica carrega consigo pressupostos sobre o que é a pessoa humana, sobre o que é o conhecimento, sobre o que é a decisão correta. Ignorar esses pressupostos é entregar à técnica a prerrogativa de definir, na prática, o conteúdo dos direitos fundamentais.

A \MH{}, ao explicitar os fundamentos antropológicos e éticos que devem orientar o desenvolvimento da IA, oferece critérios para que o direito possa exercer sua função de proteção da pessoa sem recorrer a uma rejeição simplista da tecnologia. Cinco desafios principais foram identificados:

\begin{enumerate}
    \item A \textbf{automação decisória} e seus impactos sobre a autonomia humana, que exige a afirmação do direito à explicação e à revisão humana das decisões automatizadas.
    \item A \textbf{discriminação algorítmica}, que reproduz e amplia desigualdades existentes e demanda mecanismos robustos de auditoria e não discriminação.
    \item A \textbf{vigilância e a perda de privacidade}, que criam assimetrias de poder incompatíveis com o Estado Democrático de Direito.
    \item A \textbf{fragmentação da identidade pessoal}, que impõe a necessidade de proteger a integridade da representação algorítmica da pessoa.
    \item O \textbf{transumanismo} e seus desafios aos limites do humano, que exigem uma reflexão sobre os limites éticos e jurídicos das intervenções tecnológicas na pessoa.
\end{enumerate}

O capítulo seguinte examina como o ordenamento jurídico brasileiro tem enfrentado esses desafios, identificando as potencialidades e as insuficiências dos instrumentos normativos existentes para a proteção da personalidade na era da IA.
